\section{A copy that is not a continuation}
\label{sec:copy}

We begin with a worked example. It makes the three questions concrete and
shows why endpoint similarity cannot answer them.

\paragraph{Two copies of a deadline system.}
The companion forcing paper studies a small finite system: a walker on a wheel
with a hidden phase and a deadline \citep{TsiokosForcing2026}. Its viable
states form a kernel on which a predictive self-coordinate is forced. We take
two copies of that system and tag them \emph{source} and \emph{target}. The
two tagged copies, together with the deadline system's lethal states, give a
finite state space of \snapPosZeroFourTaggedAmbientStates{} states, of which
\snapPosZeroFourTaggedLethalStates{} are lethal. To this we add the apparatus
of a handoff: a position along the episode, a budget of tokens, and a register
recording which route the episode took.

A handoff runs through four stages. In the first, only the source is live. In
the second and third, both systems are live: the source still acts, the target
already acts, and each step is supported on both. In the fourth, the source
has been withdrawn and the target continues alone. We call the middle stages
the \emph{overlap}. The target's copy of the deadline system is present from
the start, but it is latent. It becomes active only at an overlap event, with
the source present and a funded receiver. At that event the receiver reads the
position and phase at the end of the source's history and writes them into
the target. The
source is withdrawn only afterwards. This transfer rule is part of the model's
design. It does not show the target's organization being built step by step,
and it is not independent evidence about how a transfer would have to work.

On this model, conditions (1), (2), (3) and (8) of the architecture hold on
each copy in substantive form, computed from the deadline system's own
dynamics. Conditions (4) to (7) hold in the simplified form described in
Section~\ref{sec:setting}. Everything we say about this model is relative to
that grade.

\paragraph{The handoff passes.}
On this model a single handoff, from the source copy to the target copy, passes
all three checks: the target realizes the conditions at the stated grade, the
episode is a strict live-overlap lineage, and after the source and the
transfer aids are withdrawn the target keeps itself going. This is an instance
of the continuation theorem of Section~\ref{sec:continuation}, with one
difference between how lineage and autonomy draw on the source that we state
there. We now vary the picture in four ways. Each variation keeps something that
looks like success. Two fail the lineage check, one fails a typing
requirement, and one shows that public observations lose information.

\paragraph{The detached copy.}
Suppose the target's state has the same predictive readout as the
source's final state, and the target satisfies the conditions and maintains
itself, but at no stage are source and target live together: one stage has
only the source, the next only the target. We prove that each of these can
hold together with the last (Proposition~\ref{prop:detached}). The endpoint
readouts agree, the target realizes the conditions at the stated grade, and
it is autonomous. Yet no overlap ever carried the source's running
organization into it, and no strict handoff certificate exists for this
sequence of stages. The proposition is not deep. The lineage check requires a
stage with both systems live, and the point of the example is that the other
two checks, and agreement at the endpoint, can all hold without one. Whatever one wants to say about such a copy, it is
not a continuation in the sense of our lineage check, and nothing at its
endpoints reveals this. Theories that analyse persistence by psychological
continuity already treat the cause of a match as a parameter. Parfit's
versions of the psychological criterion differ exactly over which causes count
\citep[\S 78]{Parfit1984}. Continuity theorists also disagree over whether
copying mental contents supplies the right kind of causal dependence
(\citealp[\S 4]{Olson2023SEP}; compare \citealp{ShoemakerSwinburne1984} and
\citealp{Unger1990}). Our lineage check is not a criterion of persistence. Like
those debates, it looks at the route rather than at the endpoints.

\paragraph{The gap.}
Suppose instead that the episode has the right endpoints, but the stream stops
for a while and then resumes. We prove that on the two-copy model, removing one
overlap stage from the realized episode leaves the source and target stages in
place, leaves the two routes' endpoints in agreement, and keeps equal predictive
readouts at the endpoint, while the no-gap condition of the lineage check
fails
(Proposition~\ref{prop:gap}). In words, it is a resumed stream rather than a
continuation without a gap. A stage is missing from the realized episode, so
the no-gap check fails. We do not claim that anything is actively transferred
across the gap.

\paragraph{The re-encoding.}
Suppose the ``target'' is the source system itself, with its states relabelled
by a nontrivial permutation. Such a re-encoding does not meet the requirement
that source and target be distinct systems, so it is not transport between
systems, however different the relabelled description looks
(Remark~\ref{rem:reencoding}). The point is one of typing, not a theorem.

\paragraph{The public shadow.}
Finally, suppose we judge the target only by what it shows through its public
channel. We prove that two target states can agree on every observation
through that fixed channel while their predictive states differ
(Proposition~\ref{prop:shadow}). Public observations through this interface
do not determine the predictive state. Whatever a realization check needs to
know about the predictive state, it cannot read it off this interface. This is
a statement about the fixed public interface of this finite model. It is not a claim that no observation could ever reveal the
predictive state.

\paragraph{What the example shows.}
Each variation keeps something that a criterion based on endpoints or public
behaviour would count as success. The detached copy and the gap each fail the
lineage check. The re-encoding fails a typing requirement before any check
applies. The public shadow fails no check: it shows that public observations
lose information about the predictive state, which the checks use. That is why we state realization, lineage and autonomy as separate predicates,
and why the certificate in Section~\ref{sec:continuation} checks them
separately. The lower rows of Figure~\ref{fig:three} draw the detached copy and
the gap beside the passing handoff.
